% ===================================================================
%  TABELO N-RO 2 — La homa korpo. — La paŭzo. — La ludoj.
%  Traduction 2026 d'apres texto/io, controlee sur texto/fr et
%  texto/es. Consigne : voir texto/eo/10-tabelo-01.tex.
% ===================================================================

%%K t02-apar-1 apar
\VUsaut{4.0mm}
\VUtitre{15.0pt}{60}{TABELO N\textsuperscript{o} 2}
\VUsaut{2.0mm}
\VUtitre{11.4pt}{0}{La homa korpo. --- La paŭzo.}
\VUtitre{11.4pt}{0}{La ludoj.}

%%K t02-tit-0 sub
\VUtitre{13.2pt}{0}{La korpo homa.}
\VUsaut{2.4mm}
\VUcentre{10.2pt}{0}{\textit{(Simpla leciono pri naturscienco.)}}
\VUsaut{3.0mm}

%%K t02-tit-1 sub pos=t02-01-1
\VUpk{10.2pt}{40}{La kapo.}
\VUsaut{3.0mm}

%%K t02-01-1 p
1. --- La ĉefaj partoj de la homa korpo estas: la \VUgras{kapo}, la
\VUgras{torso} kaj la \VUgras{membroj}.

%%K t02-02-1 p
2. --- La kapo enhavas du partojn: la \VUgras{kranion}, kiu entenas la
\VUgras{cerbon} kaj kiun kovras la \VUgras{hararo}, kaj la \VUgras{vizaĝon}.

%%K t02-03-1 p
3. --- Sur nia desegnaĵo ni vidas nek la kranion nek la cerbon; sed ni tre
klare vidas la gravajn partojn de la vizaĝo.

%%K t02-04-1 p
4. --- Jen unue la \VUgras{frunto}, kiu montras potencan intelekton, pensadon
ĉiam agantan. La \VUgras{tempioj} estas tre liberaj. La \VUgras{okulo} estas
vigla kaj large malfermita. Densa \VUgras{brovo} kaj longaj \VUgras{okulharoj}
ŝirmas ĝin.

%%K t02-05-1 p
5. --- Jen krome la \VUgras{nazo} kaj la \VUgras{naztruoj}. La nazo servas al
ni por flari kaj por spiri. Ĉiuflanke de la nazo estas la \VUgras{vangoj},
pli malproksime la \VUgras{oreloj}. La malsupro de la vizaĝo konsistas el la
\VUgras{buŝo} kaj la \VUgras{mentono}.

%%K t02-06-1 p
6. --- Ni ne tre bone vidas ilin, ĉar la \VUgras{lipbarbo} kaj la
\VUgras{vangbarbo} kaŝas ilin de ni. Sed ni scias, ke la buŝo havas la du
\VUgras{lipojn}, la \VUgras{supran makzelon} kaj la \VUgras{malsupran
makzelon} kun iliaj \VUgras{dentoj}, kaj la \VUgras{langon}. Per la buŝo ni
parolas kaj manĝas. Per la lango ni gustumas la nutraĵojn.

%%K t02-07-1 p
7. --- La okulo, la orelo, la nazo kaj la buŝo estas, same kiel la fingroj,
la organoj de la kvin sensoj: la \VUgras{vidsenso}, la \VUgras{aŭdsenso}, la
\VUgras{flarsenso}, la \VUgras{gustsenso}, la \VUgras{tuŝsenso}.

%%K t02-08-1 p
8. --- La kapo estas do unu el la plej interesaj partoj de la homa korpo.

%%K t02-tit-2 sub
\VUsaut{4.4mm}
\VUpk{10.2pt}{40}{La torso.}
\VUsaut{3.0mm}

%%K t02-09-1 p
9. --- La \VUgras{kolo} kunigas la kapon al la torso. Ĝi enhavas la
\VUgras{gorĝon} kaj la \VUgras{nukon}.

%%K t02-10-1 p
10. --- La torso ankaŭ enhavas multajn esencajn organojn. En la \VUgras{brusto}
estas la \VUgras{pulmoj}, per kiuj ni spiras, kaj la \VUgras{koro} kiu estas
la centro de la vivo. Kiam la koro ĉesas bati, la vivanta estaĵo mortas.

%%K t02-11-1 p
11. --- La \VUgras{ripoj} tenas la bruston.

%%K t02-12-1 p
12. --- La ceteraj partoj de la torso estas la \VUgras{flanko}, la
\VUgras{dorso}, la \VUgras{ŝultroj} kaj la \VUgras{abdomeno} (ventro). Ni
kuŝiĝas sur la flanko aŭ sur la dorso por dormi, ni portas la ŝarĝojn sur nia
ŝultro.

%%K t02-13-1 p
13. --- En la abdomeno estas la \VUgras{stomako} kaj la \VUgras{intestoj}. Ili
servas al ni por digesti la nutraĵojn.

%%K t02-tit-3 sub
\VUsaut{4.4mm}
\VUpk{10.2pt}{40}{La membroj.}
\VUsaut{3.0mm}

%%K t02-14-1 p
14. --- Ni havas kvar \VUgras{membrojn}: du \VUgras{brakojn} kaj du
\VUgras{krurojn}. Per la brakoj ni prenas, tenas, levas, kolektas la objektojn;
per la kruroj ni staras, marŝas kaj kuras.

%%K t02-15-1 p
15. --- La \VUgras{ŝultro} kunigas la brakon al la korpo. Tre fortika
\VUgras{muskolo}, la biceps, movas la brakon kiun la \VUgras{kubuto} kunigas
al la \VUgras{antaŭbrako}. La \VUgras{karpo} konsistigas la artikon de la
\VUgras{mano}, kiu finiĝas per la \VUgras{fingroj} kaj la \VUgras{ungoj}. Sur
la mano ni distingas \VUgras{vejnon} (kaj arterion) en kiu fluas la
\VUgras{sango}.

%%K t02-16-1 p
16. --- La partoj de la kruro estas: la \VUgras{femuro}, la \VUgras{genuo}
kaj la \VUgras{poplito}, la \VUgras{tibikarno}, kies \VUgras{karnon} kovras la
\VUgras{haŭto}, la \VUgras{maleolo} kaj la \VUgras{piedo}. La kruraj
\VUgras{ostoj} estas tre dikaj kaj tre fortaj. Ĉe la ekstremo de la piedo
estas la \VUgras{piedfingroj}. La ceteraj partoj estas la \VUgras{plando} kaj
la \VUgras{kalkano}.

%%K t02-17-1 p
17. --- Tia estas la preskaŭ kompleta priskribo de la homa korpo.

%%K t02-17-2 p
\textit{Noto.} --- \textit{Per helpo de la esprimo: « mi doloras je... (la
kapo k. c.) » la instruisto povos reuzi ĉi tiun tutan vortaron el la
vidpunkto de la bona aŭ malbona} \VUgras{korpa stato.}

%%K t02-tit-4 sub
\VUsaut{5.0mm}
\VUpk{10.2pt}{40}{Gimnastiko.}
\VUsaut{2.4mm}
\VUcentre{10.2pt}{0}{\textit{(Rakonto de unu el la lernantoj.)}}
\VUsaut{3.0mm}

%%K t02-18-1 p
18. --- Por esti fleksebla, lerta kaj sana, ni devas ekzerci niajn krurojn kaj
niajn brakojn. Pro tio ni ludas kaj praktikas \VUgras{gimnastikon}.

%%K t02-19-1 p
19. --- Dum la semajno, ĉi tiuj du ekzercoj okazas en la korto de la liceo
(vidu la tabelon n-ro 1). Sed ĵaŭde kaj dimanĉe ni iras sur grandan
herbejon, kie ni havas spacon por kuri kaj amuziĝi.

%%K t02-20-1 p
20. --- Niaj majstroj kaj niaj gepatroj kelkfoje akompanas nin tien; sed
\VUgras{instruisto pri gimnastiko} \textsuperscript{(12)} ja direktas la
ekzercojn.

%%K t02-21-1 p
21. --- Sur la herbejo, maldekstre de la enirejo, troviĝas la
\VUgras{gimnastikaj aparatoj} \textsuperscript{(1)}, ni grimpas per la
\VUgras{ŝtupetaro} \textsuperscript{(2)}, ni renverse agas ĉe la
\VUgras{ringoj} \textsuperscript{(3)} kaj ĉe la \VUgras{trapezo}
\textsuperscript{(4)}, ni hisas nin per la manoj ĝis la supro de la
\VUgras{ŝnura ŝtupetaro} \textsuperscript{(5)} aŭ de la \VUgras{ŝnuro noda}
\textsuperscript{(6)}.

%%K t02-22-1 p
22. --- Dume, niaj kamaradoj saltas super la \VUgras{ligna ĉevalo}
\textsuperscript{(7)} aŭ la \VUgras{stangoj paralelaj} \textsuperscript{(11)}.
Aliaj \VUgras{boksas} \textsuperscript{(8)} aŭ \VUgras{skermas}
\textsuperscript{(9)} kun maskilo kaj per \VUgras{floreto}. Fine, aliaj
\VUgras{haltrolevas} \textsuperscript{(10)}.

%%K t02-tit-5 sub
\VUsaut{4.6mm}
\VUpk{10.2pt}{40}{La ludoj.}
\VUsaut{3.0mm}

%%K t02-23-1 p
23. --- Ĉirkaŭ la gimnastikantoj, knaboj kaj knabinoj ludas per diversaj
\VUgras{ludoj}. La knabinoj amuziĝas per sia \VUgras{rondego}
\textsuperscript{(14)}; ili saltas per \VUgras{ŝnuro} \textsuperscript{(13)}
aŭ invitas siajn fratojn kaj kuzojn al \VUgras{kroketa}
\textsuperscript{(15)} partio. Ili ĝojas forpuŝi la \VUgras{buleton} de la
kontraŭulo per sia \VUgras{ligna marteleto}, ĝuste kiam li pensas facile
trairi la arketon!

%%K t02-24-1 p
24. --- La knaboj preferas la ludojn pli muskolajn. Volonte ili ludas per
\VUgras{kegloj} \textsuperscript{(16)} kiujn ili faligas lerte per
\VUgras{globo} \textsuperscript{(17)}. Ili ankaŭ ne malŝatas ludi per
\VUgras{turbo} \textsuperscript{(18)} kaj \VUgras{bilboketo}
\textsuperscript{(19)} aŭ eĉ \VUgras{ranludo} \textsuperscript{(20)}. Sed
iliaj preferataj ludoj estas ja la \VUgras{buletoj} \textsuperscript{(21)} kaj
la \VUgras{murpilkado} \textsuperscript{(22)}, ĉar la muro de la
\VUgras{varmejo} \textsuperscript{(26)} estas tre taŭga por resaltigi la
malpezan pilkon.

%%K t02-25-1 p
25. --- La malgranda Johano, iranta per siaj \VUgras{stilzoj}
\textsuperscript{(24)}, estas tre lerta kaj tre bone scias konservi la
ekvilibron. Apud li, kelkaj kamaradoj ludas je \VUgras{kaŝoserĉo}
\textsuperscript{(25)} en tiu varmejo kaj malantaŭ la \VUgras{heĝo}. Aliaj
preferas la \VUgras{frapdivenon} \textsuperscript{(28)} aŭ la
\VUgras{flugpilkon} \textsuperscript{(29)} kiu flugetas de unu \VUgras{rakedo}
al la alia.

%%K t02-noto-1 noto
\VUnotes{143.2mm}{%
(*) La infanaj aŭ knabaj ludoj havas ofte nomojn pure naciajn. Ni agis por
nomi ilin en Ido kiel eble plej bone, jen inspiriĝante per la ago mem, per kiu
ili estas karakterizataj, jen elektante la nacian nomon en ĉi tiuj aŭ tiuj
landoj, kiam ĝi ne estas tro maljusta. Ekz.: la \VUgras{barelludo} (germane
kaj france) estas la \VUgras{ranludo} \textsuperscript{(20)} (angle) en kiu la
ludantoj penas ĵeti siajn rondajn diskojn tra la buŝo de la metala rano kiu
gapante staras sur la kesto (barelo) truigita. La \VUgras{ŝultrosalto}
\textsuperscript{(30)} diferencas de la \VUgras{dorsosalto} nur per tio: por
salti super la kapo, la ludantoj apogas la manojn sur la ŝultroj de sia
partnero, dum en la « \VUgras{dorsosalto} » ili apogas la manojn nur sur lia
dorso. --- Aŭ ankaŭ kelkaj el la partoprenantoj estas kliniĝintaj dum la aliaj
saltos super ilia dorso apogante la manojn sur la ŝultro; la unuaj devas
toleri la ŝokon sen subfleksiĝi, la duaj devas salti sen perdi la ekvilibron
(vidu la tabelon mem).%
}

%%K t02-26-1 p
26. --- Meze de la herbejo estas du arboj. Ĉe la piedo de unu el la du, sep aŭ
ok knaboj organizis partion de \VUgras{ŝultrosalto} \textsuperscript{(30)}
(*). Ne en ĉiuj landoj oni konas tiun ludon; sed ĉiuj scias kio estas la
\VUgras{dorsosalto}, per kiu la knaboj ne ludas ĉi-foje.

%%K t02-tit-6 sub
\VUsaut{5.0mm}
\VUpk{10.2pt}{40}{La ludoj en libera aero.}
\VUsaut{3.4mm}

%%K t02-27-1 p
27. La \VUgras{blindludo} \textsuperscript{(31)} estas ankaŭ tre amuza;
knabinoj kaj knaboj alterne metas la \VUgras{okulbendon} ĉe siaj okuloj kaj
kaptas la mallertulojn kiuj lasas sin kapti.

%%K t02-28-1 p
28. --- Meze de la herbejo, jen ludo tre franca sed iom perforta: tio estas
la \VUgras{ursino} \textsuperscript{(32)} multe pli brua ol la tiel trankvila
ludo de la \VUgras{kajto} \textsuperscript{(33)} aŭ eĉ ol la angla ludo
\VUgras{teniso} \textsuperscript{(34)}.

%%K t02-29-1 p
29. --- Tiu lasta sporto estas la ĝojo de la grandaj lernantoj. Niaj
instruistoj mem volonte praktikas ĝin. Ĝi postulas grandan lertecon kun
facilmoveco kaj ĝi tre plaĉas al mi. Mi lasas al la knabinoj la ludojn de
\VUgras{balancilo} \textsuperscript{(35)}.

%%K t02-30-1 p
30. --- Mi ne ŝatas alian anglan ludon kiu eble fariĝus danĝera.

%%K t02-30-1b p
Mi intencas paroli pri la \VUgras{piedpilko} \textsuperscript{(36)} kiu
ĝojigas mian pli aĝan fraton. Mi preferas nian malnovan \VUgras{barludon}
\textsuperscript{(38)} tiel same higienan kaj vere malpli brutalan.

%%K t02-tit-7 sub
\VUsaut{4.6mm}
\VUpk{10.2pt}{40}{Biciklado kaj pilkludo.}
\VUsaut{3.0mm}

%%K t02-31-1 p
31. --- Ankaŭ volonte mi ĉirkaŭiras la herbejon per \VUgras{biciklo}
\textsuperscript{(37)}.

%%K t02-32-1 p
32. --- Dum la venonta jaro mi lernos la \VUgras{rajdadon}
\textsuperscript{(39)}. Tiel bela estas ĉevalo kiu gracie paŝas aŭ trotas, kaj
kiu galope saltas super la obstakloj!

%%K t02-33-1 p
33. --- Somere, ni lernas la \VUgras{naĝarton} \textsuperscript{(40)}. Ni
delektiĝante plonĝas kaj banas nin en la klara kaj freŝa akvo de la rivero.
Kelkfoje fine ni lanĉas paperan \VUgras{balonon} \textsuperscript{(41)} kiu
majeste supreniras en la atmosferon, tuj kiam ĝi estas plenigita per varma
aero.

%%K t02-34-1 p
34. --- Estas ankaŭ multaj aliaj ludoj, sed mi ne finus enumeri ilin kaj oni
konscias laŭ ĉi tiu resumo, ke en la ĵaŭdoj oni ne tre enuas sur nia bela
herbejo.

%%K t02-34-2 p
N.-B. --- \textit{La instruisto facile kompletigos tiun ĉapitron; pli detale
priskribante ĉiun el la plej gravaj ludoj.}
\VUsaut{6.0mm}
\VUfilet{22mm}
